\section{A Framework for Benchmarking Open Government Data Efforts}
Der Artikel ``A Framework for Benchmarking Open Government Data Efforts''\cite{open_gov} wurde ausgewählt um zu vergleichen, wie andere Länder ihre Daten veröffentlichen. Der wissenschaftliche Artikel wurde von IEEE veröffentlicht und ist öffentlich abrufbar. 
\smallskip
\noindent
Der Artikel vergleicht wie 35 verschiedene Länder ihre Daten veröffentlicht haben. 

\smallskip
Seit 2009 haben immer mehr Länder ihre Daten veröffentlicht. Angefangen hat dieser Trend mit data.gov, die als eine Initiative der Obama Regierung geschaffen wurde. Weltweit folgten Regierungen diesem Trend, dabei auch Autokratien wie Russland und die Vereinigten Arabischen Emirate.
Laut dem Artikel bieten rund 66\% der Portale Daten in leicht Maschinen-lesbaren Formaten an. Dies erlaubt es NutzerInnen, die veröffentlichten Daten leichter zu verarbeiten und zu benutzen. Ein Großteil der Portale erlaubt NutzerInnen, mit ihnen in verschiedenen Wegen zu interagieren. Rund die Hälfte der Portale erlauben es NutzerInnen, Anfragen für Datensätze zu stellen und rund 57\% erlauben es NutzerInnen, im Portal zu diskutieren. 
Die Review nennt aber auch mehrere der Probleme, die Open Government Data zur Zeit der Publikation hatte.
So wurde in vielen Fällen nicht darauf geachtet, welche Daten für die Öffentlichkeit interessant sind, sondern darauf, welche einfach zu veröffentlichen sind. Auch wurden Daten oft in schlecht verarbeitbaren Dateiformaten publiziert, die Vereinigten Arabischen Emirate veröffentlichten viele Dateien als PDFs oder Word Dateien, die schwer maschinell auslesbar sind. Ein Grund für diese Probleme ist wohl, dass Daten mit nur wenig Vorbereitung veröffentlicht worden sind. Die Gründe dafür sind wohl meist politisch, da Länder beim von den USA mit data.gov losgetretenen "Trend" mitmachen wollten. Bei vorzeitiger Vorbereitung der Datensätze kann schon in der Planung von Open Government Data Portalen bedacht werden, wie man am besten auf sie zugreifen können sollte. Es sollte nach der Veröffentlichung kontinuierlich beachtet werden, auf welche Daten NutzerInnen zugreifen um diese priorisiert zu veröffentlichen.